\section{Variabel, Tipe Data, dan Kontrol Alur}

\begin{learningoutcomes}
\item Menggunakan tipe data dasar C\# dan tipe Unity (Vector3, Quaternion)
\item Menggunakan public, private, dan SerializeField
\item Menerapkan if-else, switch, loop dalam konteks game
\end{learningoutcomes}

\subsection{Tipe Data Dasar dan Unity}

\textbf{Dasar}: int, float, bool, string, char. \textbf{Unity}: Vector3, Vector2, Quaternion, Color, GameObject, Transform.

\begin{lstlisting}[language={[Sharp]C}]
using UnityEngine;

public class VariableExample : MonoBehaviour
{
    public int health = 100;
    public float speed = 5.5f;
    public Vector3 position;
    
    [SerializeField] private int maxHealth = 100;
}
\end{lstlisting}

\subsection{Kontrol Alur}

\begin{lstlisting}[language={[Sharp]C}]
// If-else
if (health > 0) { /* Player hidup */ }
else { GameOver(); }

// Switch
switch (gameState) {
    case GameState.Menu: ShowMenu(); break;
    case GameState.Playing: StartGame(); break;
}

// Loop
for (int i = 0; i < 10; i++) Debug.Log(i);
foreach (GameObject enemy in enemies) enemy.SetActive(false);
\end{lstlisting}

\begin{catatan}
Public variables muncul di Inspector. Gunakan SerializeField untuk private variables yang perlu diatur di Inspector.
\end{catatan}
